% source: An algebraic approach to the Riemann-Roch theorem and the arithmetic theory of function fields (Athanasios Papaioannou), http://www.fen.bilkent.edu.tr/~franz/mat/Riemann-Roch.pdf
% pdf_page: 0005
% retrieved: 2026-07-14
% transcribed_by: page-transcriber (AI vision, no OCR)

\documentclass[11pt]{article}
\usepackage{amsmath,amssymb,amsthm}

% Small-caps theorem style matching the source (DEFINITION/THEOREM/COROLLARY headings, italic bodies).
\newtheoremstyle{scplain}{\topsep}{\topsep}{\itshape}{0pt}{\scshape}{.}{ }{\thmname{#1} \thmnote{#3}}
\theoremstyle{scplain}
\newtheorem*{theorem}{Theorem}
\newtheorem*{definition}{Definition}
\newtheorem*{corollary}{Corollary}

% Small-caps "Proof." to match the source (standard amsthm definition with \scshape head).
\makeatletter
\renewenvironment{proof}[1][\proofname]{%
  \par\pushQED{\qed}\normalfont\topsep6\p@\@plus6\p@\relax
  \trivlist\item[\hskip\labelsep\scshape #1\@addpunct{.}]\ignorespaces}%
  {\popQED\endtrivlist\@endpefalse}
\makeatother
\renewcommand{\qedsymbol}{$\square$}

% This page uses only standard operators (\min) and standard amssymb symbols
% (\mathbb, \mathfrak, \mathcal, \mathbf, \varnothing, \subsetneq); no custom operators required.
% (The top line completes the proof of Theorem 1.5 begun on the preceding page.)

\begin{document}

\noindent
and this contradicts are assumption that $z_1(P), z_2(P), \ldots, z_n(P)$ are
$\mathbb{K}$-independent.\hfill$\square$

\medskip
Stretching the analogy with the geometric case, we now proceed to define the
meaning of \emph{zeroes} and \emph{poles} of elements of $\mathbb{F}/\mathbb{K}$.
We have

\begin{definition}[1.5]
Let $z \in \mathbb{F}_P$ and $P \in \mathbf{P}_{\mathbb{F}}$. We call $P$ a zero of
$z$ if and only if $v_P(z) = m > 0$; in this case $P$ is a zero of order $m$. We
call $P$ a pole of $z$ if and only if $v_P(z) = m < 0$; in this case $P$ is a pole
of order $-m$.
\end{definition}

Our next objective will be to show that any given element of $\mathbb{F}/\mathbb{K}$
has at least one pole and at least one zero. This claim is not surprising, since
intuitively $\mathbb{F}/\mathbb{K}$ is the field of meromorphic functions of the
Riemann surface whose points are the places $P \in \mathbf{P}_{\mathbb{F}}$. We
first prove the following

\begin{theorem}[1.6]
Let $\mathbb{F}/\mathbb{K}$ be a function field and $R$ a subring of $\mathbb{F}$
such that $\mathbb{K} \subseteq R \subseteq \mathbb{F}$. If $\mathfrak{a}$ is a
proper ideal of $R$, then there exists a place $P_{\mathbb{F}}$ such that
$\mathfrak{a} \subseteq P$ and $R \subseteq \mathcal{O}_P$.
\end{theorem}

\begin{proof}
We introduce the set $\mathcal{F} = \{S \mid S$ is a subring of $\mathbb{F}$ with
$R \subseteq S$ and $\mathfrak{a}S \neq S\}$---by definition, $\mathfrak{a}S$ is
the set of all finite sums of products $a_i s_i$, where $a_i \in \mathfrak{a},
s_i \in S$ and it is an ideal of $S$. We can also order $\mathcal{F}$ by
inclusion. By the conditions of the theorem, $\mathcal{F}$ is not empty
($R \in \mathcal{F}$) and a typical Zorn's Lemma argument shows that $\mathcal{F}$
has a maximal element $\mathcal{O}$, that satisfies
$R \subseteq \mathcal{O} \subseteq \mathbb{F}$ and $\mathfrak{a}S \neq \mathcal{O}$.
It's a standard fact that $\mathcal{O}$ will be a valuation ring of
$\mathbb{F}/\mathbb{K}$ (c.f. Atiyah--MacDonald, chapter 5). We repeat the proof
below.

As $\mathfrak{a} \neq 0$ and $\mathfrak{a}\mathcal{O} \neq \mathcal{O}$, we have
$\mathcal{O} \subsetneq \mathbb{F}$ and of course
$\mathfrak{a} \subseteq \mathcal{O} - \mathcal{O}^\times$. Assume that there
existed an element $z \in \mathbb{F}$ such that $z \notin \mathcal{O}$ and
$z^{-1} \notin \mathcal{O}$; then $\mathfrak{a}\mathcal{O}[z] = \mathcal{O}[z]$ and
$\mathfrak{a}\mathcal{O}[z] = \mathcal{O}[z^{-1}]$ and we may pick
$a_0, a_1, \ldots, a_n, b_0, b_1, \ldots, b_m \in \mathfrak{a}\mathcal{O}$ such
that
\[
  1 = a_n z^n + \cdots + a_1 z + a_0
\]
and
\[
  1 = b_m z^{-m} + \cdots + b_1 z^{-1} + b_0.
\]
We may assume without loss of generality that $m \leq n$ and then we choose
$m, n$ to be minimal; even under this assumption, clearly $m, n \geq 1$.
Multiplying the two relations by $1 - b_0$ and $a_n z^n$ respectively, we obtain
\[
  1 - b_0 = (1 - b_0)a_n z^n + \cdots + (1 - b_0)a_1 z + (1 - b_0)a_0
\]
and
\[
  0 = b_m a_n z^{n-m} + \cdots + b_1 a_n z^{n-1} + (b_0 - 1)a_n z^n.
\]
Adding the above equations yields
\[
  1 = c_{n-1} z^{n-1} + \cdots + c_1 z + c_0
\]
with coefficients $c_i \in \mathfrak{a}\mathcal{O}$; this is absurd by the
minimality of $n$. Therefore, $\mathcal{O}$ is a valuation ring of
$\mathbb{F}/\mathbb{K}$, as desired.
\end{proof}

From the above, we may draw the following corollaries

\begin{corollary}[1.7]
Let $\mathbb{F}/\mathbb{K}$ be a function field, $z \in \mathbb{F}$ transcendental
over $\mathbb{K}$. Then $z$ has at least one zero and one pole; in particular,
$\mathbf{P}_{\mathbb{F}} \neq \varnothing$.
\end{corollary}

\begin{proof}
Consider the ring $R = \mathbb{K}[z]$ and the ideal $\mathfrak{a} = z\mathbb{K}[z]$.
The previous theorem asserts that there is a place $P \in \mathbf{P}_{\mathbb{F}}$
with $z \in P$, so that $P$ is a zero of $z$. Moreover, $z^{-1}$ has another $0$,
$Q \in \mathbf{P}_{\mathbb{F}}$. This $Q$ will be a pole of $z$.
\end{proof}

\begin{corollary}[1.8]
The field of constants $\tilde{\mathbb{K}}$ of $\mathbb{F}/\mathbb{K}$ is a finite
algebraic extension of $\mathbb{K}$.
\end{corollary}

\begin{proof}
Since $\mathbf{P}_{\mathbb{F}} \neq \varnothing$, we may pick a place
$P \in \mathbf{P}_{\mathbb{F}}$. Since $\tilde{\mathbb{K}}$ is embedded in
$\mathbb{F}_P$ via the residue class map
$\mathcal{O}_P \longrightarrow \mathbb{F}_P$, it follows that
$[\tilde{\mathbb{K}} : K] \leq [\mathbb{F}_P : \mathbb{K}] < \infty$.
\end{proof}

The previous corollary will prove useful in the subsequent treatment of divisors,
in that it will alow us to assume that the base field is algebraically closed;
this assumption can only be made when $\tilde{\mathbb{K}}$ is a finite extension of
the base field $\mathbb{K}$ (since it implies that $\mathbb{F}/\tilde{\mathbb{K}}$
is also a function field).

\begin{center}5\end{center}

\end{document}
